% Người phụ trách: TRANG. (~1 trang). Mục "System" trong dàn ý -> Method.
% Nguồn: HYMT2_V33_FINETUNE_20261009/README_PACKAGE.md (train), KAGGLE_A1_FAST16_20261006/README_TOI_UU.md + FIX_NOTES.md
% (runtime nộp cuối), run/kout/run_config.json (gói VLSP26-FINAL-001), vlsp-pipeline README + CHANGE_NOTES (A/B/C, bản trước).
% LƯU Ý: bản NỘP là FAST16 ("legal prose blocks"), KHÁC pipeline B (gom <=8 node, newline, khớp dòng). A/B/C chỉ là ablation bản trước.
\section{Method}
\label{sec:method}

\subsection{Backbone and fine-tuning}
We use Hy-MT2-1.8B \citep{zheng2026hymt2} (dense, fp16, about 5\,GB of VRAM), which fits the 3B limit and one Kaggle T4.

Training is causal-LM cross-entropy on the assistant turn only (translation plus end-of-sequence); prompt tokens are masked, as in our training script.
Each pair uses a ``mixed'' prompt: with probability 0.5, if glossary pairs (442 terms) occur in both source and target, a terminology-hint prompt with at most eight pairs is used, otherwise the default instruction; inference always uses the default (Appendix~\ref{app:prompts}).
Stage A1 trains LoRA \citep{hu2021lora} ($r{=}64$, $\alpha{=}128$, all attention and MLP projections, 4.15\% of parameters) on 53.5k segment-level pairs (mainly article and clause segments) for one epoch at $2\times10^{-4}$.
Stage v33 \emph{warm-starts} from the A1 adapter and trains one more epoch on the 82{,}203 pairs of data v3.3 (Section~\ref{sec:data}) at the lower rate $5\times10^{-5}$ ($2\times10^{-4}$ is the rate for training LoRA from scratch); details are in Appendix~\ref{app:hyper}.

\subsection{Document pipeline}
\begin{figure}[t]
\centering
\resizebox{\columnwidth}{!}{%
\begin{tikzpicture}[node distance=4mm and 6mm,
  box/.style={draw, rounded corners, align=center, font=\footnotesize, minimum height=8mm, inner sep=3pt},
  arr/.style={-{Stealth[length=2mm]}, thick}]
  \node[box] (parse) {Parse\\nodes};
  \node[box, right=of parse] (pack) {Pack\\blocks};
  \node[box, right=of pack] (gen) {Prompt,\\generate};
  \node[box, below=7mm of gen] (val) {Parse output,\\validate};
  \node[box, left=of val] (sel) {Keep best,\\assemble};
  \node[box, below=6mm of val] (rec) {Recovery:\\one retry};
  \draw[arr] (parse) -- (pack);
  \draw[arr] (pack) -- (gen);
  \draw[arr] (gen) -- (val);
  \draw[arr] (val) -- node[above, font=\scriptsize]{ok} (sel);
  \draw[arr] (val) -- node[right, font=\scriptsize]{failure} (rec);
  \draw[arr] (rec.east) -- ++(9mm,0) |- node[right, pos=0.3, font=\scriptsize, inner sep=2pt]{changed input} (gen.east);
  \draw[arr] (rec.west) -| node[left, pos=0.75, font=\scriptsize, inner sep=2pt]{candidate} (sel.south);
\end{tikzpicture}}
\caption{Document pipeline with one recovery round. At inference the pipeline reads only the source text.}
\label{fig:pipeline}
\end{figure}

Figure~\ref{fig:pipeline} summarises the \texttt{docmt} pipeline, written as one module per step so that each step can be changed and measured alone.

\paragraph{Parse and pack.}
A rule-based parser turns a Markdown document into \emph{nodes} (headings, article and clause labels, paragraphs, table cells, signature blocks) such that concatenating the nodes reproduces the source exactly.
In an earlier version we compared three packings (Section~\ref{sec:results}): one node per request, up to 8 neighbouring nodes (at most 768 source tokens) joined with marker tokens, and the same grouping joined only by newlines.
The submitted version (``legal prose blocks'') first merges adjacent translatable prose nodes of the same scope (headings, articles, clauses, points, list items, paragraphs) into one block, and closes a block once it exceeds 768 source tokens or 8 nodes and its last node ends with a full stop.
A new article or a new chapter or section heading starts a new block, a title chain may join the article that follows it, and tables and signature blocks stay separate units.
Each block is one translation unit and its output is not required to map line by line onto the source nodes, so no line-count check is applied to prose; instead, every source span must be covered exactly once.

\paragraph{Validate and recover.}
Each output is checked for source-span coverage, preserved headings and numbered labels (treating \vi{đ} and \texttt{dd} as the same label), and unfinished generation (no end-of-sequence before the budget, or a detected repetition loop).
A unit with such a failure gets \emph{one} recovery round with a changed input.
After a lost label the prompt adds an instruction to keep the source labels; after a VI$\to$EN unit ran to its budget the earlier (legacy) budget is used; after a repetition loop the repetition penalty is 1.15 \citep{keskar2019ctrl}.
The best translation obtained is kept; a document with unresolved units is still output but flagged, and a span for which the model returns nothing keeps the source text, also flagged, so no content is silently dropped.

\paragraph{Decoding and runtime.}
Decoding is greedy in fp16.
The output budget depends on direction, $\max(256,\lceil 2.5s\rceil{+}128)$ tokens for EN$\to$VI prose and $\max(256,\lceil 1.25s\rceil{+}128)$ for VI$\to$EN prose, where $s$ is the source length; Vietnamese references need about 1.5$\times$ more tokens than English.
Requests are batched by length to reduce padding, and the batch size (8 or 16) is calibrated on each GPU.
In an earlier comparison with the same A1 model, data, budgets and batch size 8, scheduling requests by length cut the public-test wall time of the whole job from 12{,}473 s to 8{,}080 s ($-$35.2\%), with BLEU changing by $-$0.01 (EN$\to$VI) and $+$0.05 (VI$\to$EN).

\paragraph{Anti-repetition.}
A GPU-side detector checks the generated tokens every 32 decoding steps and aborts a degenerate generation early; an aborted unit counts as unfinished and goes to the recovery round above \citep{holtzman2020curious}.
The first pass uses no repetition penalty.
With the same v33 adapter and software environment (single runs), a first-pass penalty of 1.05 scored 70.28 / 47.40 BLEU and 1.0 scored 70.87 / 47.44, so we use no penalty in the first pass; the difference comes from one run per setting and has no confidence interval.
